\par\Needspace{6\baselineskip}\medskip
\begin{center}\hypertarget{native-226-article163}{}\pdfbookmark[2]{163.}{article-163}163.\end{center}
Das seit Artikel 136 dargestellte Verfahren ist vor allem für die erste Bestimmung einer bisher völlig unbekannten Bahn geeignet. Es lässt sich aber mit gleichem Erfolg anwenden, wenn eine bereits annähernd bekannte Bahn anhand von drei beliebig weit auseinanderliegenden Beobachtungen korrigiert werden soll. In diesem Fall empfehlen sich einige Änderungen. Umfassen die Beobachtungen eine sehr große heliozentrische Bewegung, so lassen sich $\dfrac{\theta''}{\theta}$ und $\theta\theta''$ nicht mehr als Näherungswerte der Größen $P$, $Q$ betrachten. Aus den annähernd bekannten Elementen lassen sich vielmehr weit genauere Werte gewinnen. Mit diesen Elementen kann man die heliozentrischen Orte auf der Bahn für die drei Beobachtungszeiten leicht berechnen. Bezeichnet man die wahren Anomalien mit $v$, $v'$, $v''$, die Radiusvektoren mit $r$, $r'$, $r''$ und den Halbparameter mit \SourceLeaf{227}$p$, so erhält man die folgenden Näherungswerte.
\[
\begin{aligned}
P&=\frac{r\sin(v'-v)}{r''\sin(v''-v')}\mathpunct{\raisebox{.45ex}{,}}\\[.6em]
Q&=\frac{4r'^4\sin\frac12(v'-v)\sin\frac12(v''-v')}{p\cos\frac12(v''-v)}\mathpunct{\cdot}
\end{aligned}
\]
\noindent Auf diese Werte stützt man also die erste Hypothese, auf geringfügig nach Belieben veränderte Werte die zweite und dritte. Es wäre nämlich unzweckmäßig, hier $P'$ und $Q'$ als neue Werte zu übernehmen (wie wir es oben getan haben), da man nicht mehr annehmen darf, dass diese Werte genauer ausfallen. Auf diese Weise lassen sich alle drei Hypothesen sehr bequem \emph{zugleich} ausarbeiten. Die vierte bildet man dann nach den Anweisungen in Artikel 120. Im Übrigen haben wir nichts dagegen einzuwenden, wenn jemand das eine oder andere der zehn in den Artikeln 124---129 dargestellten Verfahren in einem solchen Fall für beinahe ebenso zweckmäßig, wenn nicht für zweckmäßiger hält und es daher lieber anwendet.
\par\nopagebreak[4]\smallskip{\centering\rule{.23\linewidth}{.35pt}\par}
